\documentclass[xcolor=dvipsnames,10pt]{beamer}

\usepackage{../fenna-files/packages}
\usepackage{../fenna-files/commands}
\bibliography{../fenna-files/references}{}

\usepackage{multicol}

\geometry{paperwidth=140mm,paperheight=105mm}
%
% Choose how your presentation looks.
%
% For more themes, color themes and font themes, see:
% http://deic.uab.es/~iblanes/beamer_gallery/index_by_theme.html
%
\mode<presentation>
{
\usetheme{Berlin}      % or try Darmstadt, Madrid, Warsaw, ...
\definecolor{goethe}{rgb}{0,0.37,0.66}
\usecolortheme[named=goethe]{structure}
	% \usefonttheme{serif}  % or try serif, structurebold, ...
\setbeamertemplate{navigation symbols}{}
\setbeamertemplate{caption}[numbered]
\resetcounteronoverlays{exx}
}

\renewcommand{\eachwordone}{\sffamily}
\renewcommand{\eachwordtwo}{\sffamily}
\renewcommand{\eachwordthree}{\sffamily}

%gets rid of bottom navigation bars
\setbeamertemplate{footline}{}

%gets rid of navigation symbols
%\setbeamertemplate{navigation symbols}{}


\addtobeamertemplate{navigation symbols}{}{%
\usebeamerfont{footline}%
\usebeamercolor[fg]{footline}%
\hspace{1em}%
\insertframenumber/\inserttotalframenumber
}


\title{The \tsc{r}-pronoun and postposition \tit{waar-mee} in Dutch}
\author{Fenna Bergsma}
\date{Syntax colloquium\\ \today}
\institute{Goethe-Universität Frankfurt}



\begin{document}

\begin{frame}
\titlepage

\centering{
\includegraphics[width=0.4\textwidth]{dfg-logo}
\hspace{2cm}
\includegraphics[width=0.3\textwidth]{goethe-logo}}
\end{frame}

% \section{Introduction}

\begin{frame}

\exg. Jij schildert \tbf{met} \tbf{een} \tbf{kwast}.\\
 you paint with a brush\\
 `You are painting with a brush.'\label{ex:metwat}

\pause

\exg. \tbf{Wat} zie jij?\\
 what see you\\
 `What do you see?'\label{ex:wat}

\pause

\exg. *\tbf{Met} \tbf{wat} schilder jij?\\
 with what paint you\\
 `What are you painting with?'\label{ex:metwat}

\pause

\exg. \tbf{Waarmee} schilder jij?\\
 {where with} paint you with\\
 `What are you painting with?'\label{ex:waar-mee}

\end{frame}


\begin{frame}

\exg. Womit malst du?\\
 {where with} paint you\\
 `What are you painting with?'

\pause

\exg. Waa\tbf{r}mee schilder jij?\\
 {where with} paint you with\\
 `What are you painting with?'\label{ex:waar-mee}

\end{frame}




\begin{frame}

\exg. Worauf kletterst du?\\
 {where on} climb you\\
 `What are you climbing onto?'

\exg. Wo bist du \tbf{(d)r}auf geklettert?\\
 where are you on climbed\\
 `What did you climb on?'

\pause

\exg. Waa\tbf{r} ben jij op geklommen?\\
 where are you on climbed\\
 `What did you climb on?'


\end{frame}


\begin{frame}

\exg. *Ik schilder \tbf{met} \tbf{t}.\\
 I paint with it\\
 `I am painting with it.'

\pause

\exg. Ik schilder \tbf{r} \tbf{-mee}.\\
 I paint \tsc{r} -with\\
 `I am painting with it.'\label{ex:jarmee}

\end{frame}


\begin{frame}

\centering{\Large{Why *\tit{met wat}?}}

\vspace{3em} \citep{riemsdijk1978}


\end{frame}


\begin{frame}

\centering{
*\tit{met wat} \hspace{1em} | \hspace{1em} \tit{waar-mee}}

\vspace{1em}

\pause

\begin{itemize}
\item positions are reversed \pause
\item \tit{met} becomes \tit{mee} \pause
\item \tit{wat} becomes \tit{waar}, which is syncretic with locative \citep{koopman1994}
\end{itemize}

\vspace{1em}\pause

\tsc{r}-pronouns are the result of regular spellout mechanisms

\pause

as long as all features form a proper constituent

\end{frame}


\begin{frame}

\ex.\label{ex:headed}
\ag. Ik koop de kwast \tbf{waar} \tbf{-mee} jij schildert.\\
 I buy the brush where with you paint\\
 `I buy the brush that you are painting with.'
\bg. *Ik koop de kwast \tbf{met} \tbf{wat} jij schildert.\\
 I koop the brush with what you paint\\
 `I buy the brush that you are painting with.'

\end{frame}


\begin{frame}

\ex.
\ag. \tbf{Waar} \tbf{-mee} schilder jij?\\
 where with paint you with\\
 `What are you painting with?'\label{ex:const2}
\bg. Ik koop de kwast [\tbf{waar} \tbf{-mee} jij schildert].\\
 I buy the brush where with you paint\\
 `I koop the brush that you are painting with.'\label{ex:const3}

\pause

\ex. [[jij] [[schilder] [\tbf{waar-mee}]]]\label{ex:const2stage}


\end{frame}



\begin{frame}

\ex.
\ag. Ik heb gekocht \tbf{waar} \tbf{-mee} jij schildert.\\
 I have bought where with you paint\\
 `I bought what you are painting with.'\label{ex:gekochtwaar-mee}
\bg. *Ik heb gekocht \tbf{met} \tbf{wat} jij schildert.\\
 I have bought with what you paint\\
 `I bought what you are painting with.'\label{ex:gekochtmetwat}

\pause

\ex.\label{ex:schildermet}
\ag. *Ik schilder \tbf{waar} \tbf{-mee} jij hebt gekocht.\\
 I paint where with you have bought\\
 `I paint with what you bought.'\label{ex:schilderwaar-mee}
\bg. Ik schilder \tbf{met} \tbf{wat} jij hebt gekocht.\\
 I paint with what you have bought\\
 `I paint with what you bought.'\label{ex:schildermetwat}


\end{frame}


\begin{frame}

\ex.
\ag. Ik koop \tbf{het} \tbf{schilderij}.\\
 I buy the painting\\
 `I am buying the painting.'\label{ex:kopen}
\bg. Ik schilder \tbf{met} \tbf{een} \tbf{kwast}.\\
 I paint with a brush\\
 `I am painting with a brush.'\label{ex:schilderen}

 \pause

\exg. Ik heb gekocht [\tbf{waar} \tbf{-mee} jij schildert].\\
 I have bought where with you paint\\
 `I bought what you are painting with.'\label{ex:mismatchwaar-mee}

 \pause

\exg. Ik schilder \tbf{met} [\tbf{wat} jij hebt gekocht].\\
 I paint with what you have bought\\
 `I paint with what you bought.'\label{ex:mismatchmetwat}

\end{frame}



\begin{frame}

\exg. [\tbf{Met} [\tbf{wat} [voor [potloden]]] teken jij?\\
 with what for pencils draw you\\
 `What kind of pencils are you drawing with?'\label{ex:watwasfur}

\pause

\exg. Ik wil graag thee [\tbf{met} [\tbf{wat} [suiker]]].\\
 I want please tea with some sugar\\
 `I would like to have tea with some sugar.'\label{ex:watindef}

\pause

\ex.\label{ex:summaryconst}
\a. [[met] [wat]] → [waar-mee]\label{ex:waar-meefr}
\b. [met [wat [X]]]\label{ex:metwatx}
\b. [met [[wat] [X]]]\label{ex:metwatfr}

\end{frame}


\begin{frame}

\ex. w -aa \tbf{-r} \tbf{-mee}

\ex. \tbf{met} w -a \tbf{-t}

\pause

\begin{itemize}
\item late insertion (phonology after syntax)\pause
\item individual features (no feature bundles)
\end{itemize}


\end{frame}

\begin{frame}

\ex. \begin{forest}
[\tsc{wP}
    [W, roof]
]
{\draw (.east) node[right]{⇔ \tit{w}}; }
\end{forest}\label{ex:entryw}

\citep{hachem2015}

\pause


\ex. \begin{forest}
[\tsc{deixP}
    [\tsc{deix}, roof]
]
{\draw (.east) node[right]{⇔ \tit{a}}; }
\end{forest}\label{ex:entrya}

\citep[cf.][]{lander2016}

\end{frame}


\begin{frame}

\ex.\label{ex:tsubobj}
\ag. T staat in de hal.\\
 \tsc{3sg.n.nom} stands in the hallway\\
 `It is standing in the hallway.'\label{ex:tnoclitic}
\bg. Ik zie t.\\
 I see \tsc{3sg.n.acc}\\
 `I see it.'
\bg. Ik heb t een klap gegeven.\\
 I have \tsc{3sg.n.dat} a hit given\\
 `I gave it a hit.'

\end{frame}

\begin{frame}

\label{ex:casetree}
\ex. \begin{forest} boom
	[\tsc{insP}
			[\tsc{f4}]
			[\tsc{datP}
					[\tsc{f3}]
					[\tsc{accP}
							[\tsc{f2}]
							[\tsc{nomP}
									[\tsc{f1}]
									[X
											[..,roof]
									]
							]
					]
			]
	]
\end{forest}

\citep{caha2009}

\end{frame}



\begin{frame}

\ex.
\ag. *Hij en t staan in de hoek.\\
 he and it stand in the corner\\
 `He and it are standing in the corner.'\label{ex:tcoordinated}
\bg. Hij en dit/dat staan in de hoek.\\
 he and this/that stand in the corner\\
 `He and it are standing in the corner.'\label{ex:datcoordinated}

\pause

\exg. T staat in de hal.\\
 \tsc{3sg.n.nom} stands in the hallway\\
 `It is standing in the hallway.'\label{ex:tnoclitic}

\end{frame}



\begin{frame}

\label{ex:strengthtree}
\ex. \begin{forest} boom
			[CP
					[C]
					[ΣP
							[Σ]
							[IP
									[I]
									[X
											[..,roof]
									]
							]
					]
			]
\end{forest}

\citep{cardinaletti1996}

\end{frame}


\begin{frame}

\ex. \begin{forest} boom
 [\tsc{datP}
     [\tsc{f3}]
     [\tsc{accP}
         [\tsc{f2}]
         [\tsc{nomP}
             [\tsc{f1}]
             [ΣP
                 [Σ]
								 [IP
								 		[I]
		                 [\tsc{thingP}
		                     [\tsc{thing}, roof]
										]
                 ]
             ]
         ]
     ]
 ]
 {\draw (.east) node[right]{⇔ \tit{t}}; }
 \end{forest}\label{ex:entryt}


\citep{cardinaletti1996,kayne2005,caha2009}

\end{frame}


\begin{frame}

\ex. \tbf{The Superset Principle} \citep{starke2009}: \\
	A lexically stored tree matches a syntactic node iff the lexically stored tree contains the syntactic node.

\begin{multicols}{2}

\ex. \begin{forest} boom
	[\tsc{accP}
			[\tsc{f2}]
			[\tsc{nomP}
					[\tsc{f1}]
					[ΣP
							[Σ]
							[IP
									[I]
									[\tsc{thingP}
											[\tsc{thing}, roof]
									]
							]
					]
			]
	]
\end{forest}\label{ex:tacc}

\ex. \begin{forest} boom
	[\tsc{nomP}
			[\tsc{f1}]
			[ΣP
					[Σ]
					[IP
							[I]
							[\tsc{thingP}
									[\tsc{thing}, roof]
							]
					]
			]
	]
\end{forest}\label{ex:tnom}

\end{multicols}

\end{frame}




\begin{frame}

\exg. Ik ben r al geweest.\\
 I am there already been\\
 `I have already been there.'

\pause

\ex. \begin{forest} boom
[\tsc{loc}P
   [\tsc{location}]
   [\tsc{thingP}
       [\tsc{thing}, roof]
   ]
]
{\draw (.east) node[right]{⇔ \tit{r}}; }
\end{forest}\label{ex:entryr}

\citep{baunaz2018}

\pause

\ex. \tbf{The Elsewhere Condition} (\citealt{kiparsky1973}, formulated as in \citealt{caha2020}):\\
When two entries can spell out a given node, the more specific entry wins. Under the Superset Principle governed insertion, the more specific entry is the one which has fewer unused features.

\pause

\ex. \begin{forest} boom
 [\tsc{thingP}
     [\tsc{thing}, roof]
 ]
{\draw (.east) node[right]{⇒ \tit{r}}; }
\end{forest}

\end{frame}


\begin{frame}

\centering{
\begin{tabular}{cc}
\toprule
						& \tsc{n.sg} \\
\midrule
\tsc{nom} & t         \\
\tsc{acc} & t         \\
\tsc{dat} & t         \\
\tsc{ins} & r -mee    \\
\bottomrule
\end{tabular}
}

\end{frame}


\begin{frame}

\ex. \begin{forest} boom
[\tsc{insP}
    [\tsc{f4}]
    [\tsc{datP}
        [\tsc{f3}]
        [\tsc{accP}
            [\tsc{f2}]
            [\tsc{nomP}
                [\tsc{f1}]
                [ΣP
                    [Σ]
										[IP
												[I]
										]
                ]
            ]
        ]
    ]
]
{\draw (.east) node[right]{⇔ \tit{-mee}}; }
\end{forest}\label{ex:entrymee}

\pause


\ex. \begin{forest} boom
[\tsc{insP}
    [\tsc{f4},roof]
]
{\draw (.east) node[right]{⇔ \tit{met}}; }
\end{forest}\label{ex:entrymet}

\end{frame}


\begin{frame}

[W [\tsc{deix} [\tsc{f4} [\tsc{f3} [\tsc{f2} [\tsc{f1} [Σ [I [\tsc{thing}] ] ] ] ] ] ] ] ]

\end{frame}


\begin{frame}

\ex. Merge F and \label{ex:spellout}
	 \a. Spell out FP.
	 \b. If (a) fails, attempt movement of the spec of the complement of \tsc{f}, and retry (a).
	 \b. If (b) fails, move the complement of \tsc{f}, and retry (a).

\pause

\ex. \tbf{Cyclic Override} \citep{starke2018}:\\
	Lexicalisation at a node XP overrides any previous match at a phrase contained in XP.

\pause

\ex. \tbf{Backtracking} \citep{starke2018}:\\
	When spellout fails, go back to the previous cycle, and try the next option for that cycle.\label{ex:backtracking}

\pause

\ex. \tbf{Spec Formation} \citep{starke2018}:\\
	If Merge F has failed to spell out (even after backtracking), try to spawn a new derivation providing the feature F and merge that with the current derivation, projecting the feature F at the top node.\label{ex:specformation}

\end{frame}


\begin{frame}

\ex. \begin{forest} boom
 [\tsc{thingP}
     [\tsc{thing}, roof]
 ]
{\draw (.east) node[right]{⇒ \tit{r}}; }
\end{forest}\label{ex:thingspellout}

\pause

\ex. \begin{forest} boom
[IP
   [I]
   [\tsc{thingP}
       [\tsc{thing}, roof]
   ]
]
{\draw (.east) node[right]{⇒ \tit{t}}; }
\end{forest}\label{ex:thingf1}

\end{frame}


\begin{frame}

\ex. \begin{forest} boom
	   [ΣP
	       [Σ]
				 [IP
				 		[I]
			       [\tsc{thingP}
			           [\tsc{thing}, roof]
			       ]
					]
	   ]
	{\draw (.east) node[right]{⇒ \tit{t}}; }
\end{forest}

\end{frame}


\begin{frame}

\ex. \begin{forest} boom
	[\tsc{nomP}
	   [\tsc{f1}]
	   [ΣP
	       [Σ]
				 [IP
				 		[I]
			       [\tsc{thingP}
			           [\tsc{thing}, roof]
			       ]
					]
	   ]
	]
	{\draw (.east) node[right]{⇒ \tit{t}}; }
\end{forest}

\end{frame}


\begin{frame}

\ex. \begin{forest} boom
[\tsc{datP}
    [\tsc{f3}]
    [\tsc{accP}
       [\tsc{f2}]
       [\tsc{nomP}
           [\tsc{f1}]
           [ΣP
               [Σ]
							 [IP
							 		[I]
		               [\tsc{thingP}
		                   [\tsc{thing}, roof]
		               ]
							 ]
           ]
       ]
    ]
]
{\draw (.east) node[right]{⇒ \tit{t}}; }
\end{forest}

\end{frame}




\begin{frame}

\ex. \begin{forest} boom
[\tsc{insP}
    [\tsc{f4}]
    [\tsc{datP}
        [\tsc{f3}]
        [\tsc{accP}
           [\tsc{f2}]
           [\tsc{nomP}
               [\tsc{f1}]
               [ΣP
                   [Σ]
									 [IP
									 		[I]
		                   [\tsc{thingP}
		                       [\tsc{thing}, roof]
		                   ]
										]
               ]
           ]
        ]
    ]
]
{\draw (.east) node[right]{⇒ }; }
\end{forest}\label{ex:f4no}

\end{frame}


\begin{frame}

\ex. \begin{forest} boom
	[\tsc{datP}
	    [\tsc{datP}
	        [\tsc{f3}]
	        [\tsc{accP}
	           [\tsc{f2}]
	           [\tsc{nomP}
	               [\tsc{f1}]
	               [ΣP
	                   [Σ]
										 [IP
										 		[I]
			                   [\tsc{thingP}
			                       [\tsc{thing}, roof]
												 ]
	                   ]
	               ]
	           ]
	        ]
	    ]
	    {\draw (.east) node[right]{⇒ \tit{'t}}; }
	    [\tsc{insP}
	        [\tsc{f4}]
	    ]
	    {\draw (.east) node[right]{⇒ }; }
	]
\end{forest}

\end{frame}


\begin{frame}

\ex. \begin{forest} boom
[IP
	 [\tsc{thingP}
			 [\tsc{thing}, roof]
	 ]
	 {\draw (.east) node[right]{⇒ \tit{r}}; }
	 [IP
			 [I]
	 ]
	 {\draw (.east) node[right]{⇒ \tit{-mee}}; }
]
\end{forest}

\pause

\citep{abels2003diss}

\end{frame}


\begin{frame}

\ex. \begin{forest} boom
[ΣP
    [Σ]
    [IP
       [\tsc{thingP}
           [\tsc{thing}, roof]
       ]
       [IP
           [I]
       ]
    ]
]
{\draw (.east) node[right]{⇒}; }
\end{forest}\label{ex:f1again}

\end{frame}


\begin{frame}

\ex. \begin{forest} boom
[ΣP
   [\tsc{thingP}
       [\tsc{thing}, roof]
   ]
   {\draw (.east) node[right]{⇒ \tit{r}}; }
   [ΣP
       [Σ]
       [IP
           [I]
       ]
   ]
   {\draw (.east) node[right]{⇒ \tit{-mee}}; }
 ]
\end{forest}\label{ex:f1spec}

\end{frame}


\begin{frame}

\ex. \begin{forest} boom
[\tsc{insP}
    [\tsc{thingP}
       [\tsc{thing}, roof]
    ]
    {\draw (.east) node[right]{⇒ \tit{r}}; }
    [\tsc{insP}
       [\tsc{f4}]
       [\tsc{datP}
           [\tsc{f3}]
           [\tsc{accP}
               [\tsc{f2}]
               [\tsc{nomP}
                   [\tsc{f1}]
									 [ΣP
									 		[Σ]
		                   [IP
		                       [I]
											]
                   ]
               ]
           ]
       ]
    ]
    {\draw (.east) node[right]{⇒ \tit{-mee}}; }
]
\end{forest}\label{ex:spelloutrmee}

\end{frame}


\begin{frame}

\ex. \begin{forest} boom
[\tsc{\tsc{wP}}
		[\tsc{\tsc{wP}}
				[\tsc{w}, roof]
		]
		{\draw (.east) node[right]{⇒ \tit{w}}; }
		[\tsc{deixP}
				[\tsc{deixP}
						[\tsc{deix}, roof]
				]
				{\draw (.east) node[right]{⇒ \tit{a}}; }
				[\tsc{insP}
						[\tsc{thingP}
							 [\tsc{thing}, roof]
						]
						{\draw (.east) node[right]{⇒ \tit{r}}; }
						[\tsc{insP}
							 [\tsc{f4}]
							 [\tsc{datP}
									 [\tsc{f3}]
									 [\tsc{accP}
											 [\tsc{f2}]
											 [\tsc{nomP}
													 [\tsc{f1}]
													 [ΣP
													 		 [Σ]
															 [IP
																	 [I]
															 ]
													 ]
											 ]
									 ]
							 ]
						]
						{\draw (.east) node[right]{⇒ \tit{-mee}}; }
				]
		]
]
\end{forest}\label{ex:spelloutwaar-mee}

\end{frame}



\begin{frame}

\ex. \begin{forest} boom
[CP
    [\tsc{\tsc{wP}}
        [\tsc{\tsc{wP}}
            [\tsc{w}, roof]
        ]
        {\draw (.east) node[right]{⇒ \tit{w}}; }
        [\tsc{deixP}
            [\tsc{deixP}
                [\tsc{deix}, roof]
            ]
            {\draw (.east) node[right]{⇒ \tit{a}}; }
            [\tsc{datP}
                [\tsc{f3}]
                [\tsc{accP}
                   [\tsc{f2}]
                   [\tsc{nomP}
                       [\tsc{f1}]
                       [ΣP
                           [Σ]
													 [IP
													 		[I]
		                           [\tsc{thingP}
		                               [\tsc{thing}, roof]
		                           ]
														]
                       ]
                   ]
                ]
            ]
            {\draw (.east) node[right]{⇒ \tit{t}}; }
        ]
    ]
    [CP
        [..,roof]
    ]
]
\end{forest}

\end{frame}


\begin{frame}

\begin{itemize}
	\item \tsc{r}-pronouns are nothing special, but just a result of regular spellout \pause
	\item \tit{r} is ambiguous between location and (caseless) thing \pause
	\item \tit{met} and \tit{mee} are separate lexical entries
		\begin{itemize}
			\item so their phonological realization can differ
			\item so they can differ between between being pre and post
		\end{itemize}
\end{itemize}

\end{frame}


\begin{frame}

\exg. Ik zie \tbf{t}.\\
 I see it\\
 `I see it.'\label{ex:inaniobj}

\exg. *Ik schilder \tbf{met} \tbf{t}.\\
 I paint with it\\
 `I am painting with it.'\label{ex:prephet}

\exg. Ik schilder \tbf{r} \tbf{-mee}.\\
 I paint where -with\\
 `I am painting with it.'\label{ex:preper}

\pause

\exg. Ik zie \tbf{haar/hem}.\\
 I see her/him\\
 `I see her/him.'\label{ex:aniobj}

\exg. Ik schilder samen \tbf{met} \tbf{haar/hem}.\\
I paint together with her/him\\
`I am painting together with her/him.'\label{ex:prepani}

\pause

\ex. [\tsc{f4} [\tsc{f3} [\tsc{f2} [\tsc{f1} [\tbf{phi} [Σ [I [\tsc{thing}] ] ] ] ] ] ] ]

\end{frame}



\begin{frame}

\exg. Ik heb gekocht waar mee jij schildert.\\
 I have bought where with you paint\\
 `I bought what you are painting with.'\label{ex:meestranded}

\pause

\exg. Ik heb gekocht waar jij mee schildert.\\
 I have bought where you paint\\
 `I bought what you are painting with.'\label{ex:meestranded}

\pause

\ex. \begin{forest} boom
[\tsc{insP}
    [\tsc{thingP}
       [\tsc{thing}, roof]
    ]
    {\draw (.east) node[right]{⇒ \tit{r}}; }
    [\tsc{\tbf{insP}}
       [\tsc{f4}]
       [\tsc{datP}
           [\tsc{f3}]
           [\tsc{accP}
               [\tsc{f2}]
               [\tsc{nomP}
                   [\tsc{f1}]
									 [ΣP
									 		[Σ]
			                   [IP
			                       [I]
			                   ]
			               ]
								]
           ]
       ]
    ]
    {\draw (.east) node[right]{⇒ \tit{-mee}}; }
]
\end{forest}

\citep{noonan2017dutch}

\end{frame}


\begin{frame}

\begin{tabular}{ l | l | l }
\textsc{prep} \textit{what} & \textit{where} \tsc{post} & translation \\
\hline
 aan wat & waar\textbf{aan} & `on whom/what' \\
 achter wat & waar\textbf{achter} & `behind whom/what' \\
 bij wat & waar\textbf{bij} & `at whom/what' \\
 door wat & waar\textbf{door} & `by whom/what' \\
 langs wat & waar\textbf{langs} & `alongside whom/what \\
 met wat & waar\underline{\textbf{mee}} & `with whom/what' \\
 naar wat & waar\textbf{naar} & `to whom/what' \\
 op wat & waar\textbf{op} & `on whom/what' \\
 over wat & waar\textbf{over} & `about whom/what' \\
 tot wat & waar\underline{\textbf{toe}} & `to whom/what' \\
 uit wat & waar\textbf{uit} & `out of whom/what' \\
 van wat & waar\textbf{van} & `of whom/what' \\
 voor wat & waar\textbf{voor} & `for whom/what' \\
\end{tabular}

\end{frame}


\begin{frame}

\ex.\label{ex:dutchin}
\ag. Ik dans in het bos.\\
 I dance in the forest\\
 `I am dancing in the forest.'
\bg. Ik dans het bos in.\\
 I dance the forest in\\
 `I am dancing into the forest.'

\pause

\exg. Ik dans r -in.\\
 I dance \tsc{r} in\\
 `I dance it in.'\\
  *`I dance into it.'

\pause

\exg. *Ik ben r -in gedanst.\\
 I am \tsc{r} in danced\\
 `I dancing into it.'

\pause

\exg. Ik heb r -in gedanst.\\
 I have \tsc{r} in danced\\
 `I dancing in it.'


\end{frame}

\begin{frame}

\centering{
	\begin{tabular}{ccc}
	\toprule
					& DP				& \tsc{r}	\\
	\midrule
		pre		& \tsc{loc} & *     	\\
		post	& \tsc{dir}	& \tsc{loc} \\
	\bottomrule
\end{tabular}
		}

\vspace{1em}
\pause

verbal particles in e.g. \tbf{mee}lopen `walk along', \tbf{mee}zingen `sing along'

\vspace{1em}
	\pause

\centering{
	\begin{tabular}{ccc}
	\toprule
					& \tit{mee}	& \tit{met}	\\
	\midrule
		pre		& verb		 	& noun     	\\
		post	& noun			& *					 \\
	\bottomrule
	\end{tabular}
	}


\end{frame}




\begin{frame}{The data}

\exg. Zij schrijft de brief \tbf{met} een potlood.\\
 she writes the letter with a pencil\\
 `She writes the letter with a pencil.'

\pause

\exg. Zij koopt \tbf{dat}.\\
 she buys that\\
 `She buys that.'

\pause

\exg. * Zij schrijft de brief \tbf{met} \tbf{dat}.\\
 she writes the letter with that\\
 `She writes the letter with that.'

\pause

\exg. Zij schrijft de brief \tbf{daar} \tbf{mee}.\\
 she writes the letter there with\\
 `She writes the letter with that.'

\end{frame}



\begin{frame}{Questions and answers}

\centering{
\tbf{met dat} | \tbf{daar mee} \pause
}

\vspace{1em}

\begin{enumerate}
  \item Q: why *\tit{met dat} (but \tit{daar mee})? \pause\\
  A: \tit{daar mee} is the default, \tit{met dat} only appears under special circumstances (morphological derivation, cf. \tit{mice} and not \tit{mouse-s})\pause
  \item Q: why does the adposition change form? \pause\\
  A: \tit{mee} and \tit{met} are stored as different forms in the lexicon \pause
  \item Q: why does the adposition change position? \pause\\
  A: \tit{mee} and \tit{met} are stored as preposition or postposition in the lexicon \pause
  \item Q: why is there an \tit{r}, even though there is no location involved? \pause\\
  A: \tit{r} is ambiguous between \tsc{thing} and \tsc{location}
\end{enumerate}

\end{frame}


\begin{frame}{Nanosyntax}

\begin{itemize}
  \item functional sequence is constant\pause
  \item spellout algorithm is constant\pause
  \item a single place for cross-linguistic differences: the postsyntactic lexicon\pause
\end{itemize}

\vspace{1em}

\begin{enumerate}
  \item spellout happens feature by feature and targets phrases
  \item phonological forms are stored in the postsyntactic lexicon
  \item postpositions are candidates for spellout before prepositions are
  \item ambiguity comes from containment relations
\end{enumerate}

\end{frame}



\begin{frame}{differing phonological forms}

\begin{itemize}
  \item \tit{mee} and \tit{met} are stored as different forms in the lexicon
  \item phonological forms are stored in the lexicon \pause
\end{itemize}

  \ex. \begin{forest} boom
  [\tsc{insP}
      [\tsc{f4}]
      [\tsc{datP}
          [\tsc{f3}]
          [\tsc{accP}
              [\tsc{f2}]
              [\tsc{nomP}
                  [\tsc{f1}]
                  [ΣP
                      [Σ]
  										[IP
  												[I]
  										]
                  ]
              ]
          ]
      ]
  ]
  {\draw (.east) node[right]{⇔ \tit{-mee}}; }
  \end{forest}\label{ex:entrymee}

  \pause

  \ex. \begin{forest} boom
  [\tsc{insP}
      [\tsc{f4},roof]
  ]
  {\draw (.east) node[right]{⇔ \tit{met}}; }
  \end{forest}\label{ex:entrymet}

\end{frame}




\begin{frame}{pre vs. post}

\begin{itemize}
  \item \tit{mee} and \tit{met} are stored as preposition or postposition in the lexicon \pause
  \item postpositions are candidates for spellout before prepositions are \pause
\end{itemize}

  \ex. \begin{forest} boom
  [\tsc{insP}
      [\tsc{f4}]
      [\tsc{datP}
          [\tsc{f3}]
          [\tsc{accP}
              [\tsc{f2}]
              [\tsc{nomP}
                  [\tsc{f1}]
                  [ΣP
                      [Σ]
  										[IP
  												[I]
  										]
                  ]
              ]
          ]
      ]
  ]
  {\draw (.east) node[right]{⇔ \tit{-mee}}; }
  \end{forest}\label{ex:entrymee}

  \pause

  \ex. \begin{forest} boom
  [\tsc{insP}
      [\tsc{f4},roof]
  ]
  {\draw (.east) node[right]{⇔ \tit{met}}; }
  \end{forest}\label{ex:entrymet}

\end{frame}



\begin{frame}{pre vs. post}

  \ex. Merge F and \label{ex:spellout}
  	 \a. Spell out FP.
  	 \b. If (a) fails, attempt movement of the spec of the complement of \tsc{f}, and retry (a).
  	 \b. If (b) fails, move the complement of \tsc{f}, and retry (a).

  \ex. \tbf{Backtracking} \citep{starke2018}:\\
  	When spellout fails, go back to the previous cycle, and try the next option for that cycle.\label{ex:backtracking}

  \ex. \tbf{Spec Formation} \citep{starke2018}:\\
  	If Merge F has failed to spell out (even after backtracking), try to spawn a new derivation providing the feature F and merge that with the current derivation, projecting the feature F at the top node.\label{ex:specformation}

\end{frame}


\begin{frame}{\tit{r} without location}

  \begin{itemize}
    \item \tit{r} is ambiguous between \tsc{thing} and \tsc{location}
    \item ambiguity comes from containment relations
  \end{itemize}

\pause

\ex. \tbf{The Superset Principle} \citep{starke2009}: \\
	A lexically stored tree matches a syntactic node iff the lexically stored tree contains the syntactic node.

\ex. \begin{forest} boom
[\tsc{loc}P
   [\tsc{location}]
   [\tsc{thingP}
       [\tsc{thing}, roof]
   ]
]
{\draw (.east) node[right]{⇔ \tit{r}}; }
\end{forest}\label{ex:entryr}

\end{frame}


\begin{frame}{Derivational system}

\begin{itemize}
  \item \tit{daar mee} is the default, \tit{met dat} only appears under special circumstances (morphological derivation, cf. \tit{mice} and not \tit{mouse-s}) \pause
  \item spellout happens feature by feature and targets constituents
\end{itemize}



\pause

\ex.
\ag. Ze heeft de brief \tbf{daar} \tbf{mee} geschreven.\\
 she has the letter there with written\\
 `She wrote the letter with that.'
\bg. Zij heeft de brief \tbf{met} \tbf{dat} potlood geschreven.\\
 she has the letter with that pencil written\\
 `She wrote the letter with that potlood.'

\end{frame}


\begin{frame}{\tit{daar mee} in a proper constituent}

\exg. Ze heeft de brief \tbf{daar} \tbf{mee} geschreven.\\
 she has the letter there with written\\
 `She wrote (the letter) with that.'

\pause

\exg. \tbf{Daar} \tbf{mee} heeft ze de brief geschreven.\\
 there with has she the letter written\\
 `With that she has written the letter.' \hfill topicalization

\pause

\exg. Ze heeft de brief \tbf{hier} \tbf{mee} en \tbf{daar} \tbf{mee} geschreven.\\
 she has the letter here with and there with written\\
 `She has written the letter with this and with that' \hfill coordination

\pause

\ex.
\a. With what did she write the letter?
\bg. \tbf{daar} \tbf{mee}\\
 there with\\
 `with that' \hfill fragment answer

\end{frame}



\begin{frame}

  \ex. \begin{forest} boom
  [\tsc{insP}
      [\tsc{f4}]
  	    [\tsc{datP}
  	        [\tsc{f3}]
  	        [\tsc{accP}
  	           [\tsc{f2}]
  	           [\tsc{nomP}
  	               [\tsc{f1}]
  	               [ΣP
  	                   [Σ]
  										 [IP
  										 		[I]
  			                   [\tsc{thingP}
  			                       [\tsc{thing}, roof]
  												 ]
  	                   ]
  	               ]
  	           ]
  	        ]
  	    ]
  	    {\draw (.east) node[right]{⇒ \tit{'t}}; }
    ]
  \end{forest}


\end{frame}






\begin{frame}{\tit{met dat} in a non-constituent}


\exg. Ze heeft de brief \tbf{met} \tbf{dat} potlood geschreven.\\
she has the letter with the pencil written\\
`She has written the letter with that pencil.'

\pause

\ex.
\ag. *\tbf{Met} \tbf{dat} heeft ze de brief \tbf{potlood} geschreven.\\
 with that has she the letter pencil written\\
 `With that pencil she has written the letter.'
\bg. \tbf{Met} \tbf{dat} \tbf{potlood} heeft ze de brief geschreven.\\
 with that pencil has she the letter written\\
 `With that pencil she has written the letter.' \hfill topicalization

\pause

% \ex.
% \ag. *Ze heeft de brief met die en met dat potlood geschreven.\\
%  she has the letter with this and with that pencil written\\
%  `She has written the letter with this pencil and with that marker.' \hfill coordination
% \bg. Ze heeft de brief met die pen en en dat potlood geschreven.\\
%  she has the letter with this marker and with that pencil written\\
%  `She has written the letter with this pencil and with that marker.' \hfill coordination
%
% \pause

\ex.
\a. With what did she write the letter?
\bg. *\tbf{met} \tbf{dat}\\
 with that\\
 `with that'
\bg. \tbf{met} \tbf{dat} \tbf{potlood}\\
 with that pancil\\
 `with that pencil'\hfill fragment answer

\end{frame}

\begin{frame}

  \ex. \begin{forest} boom
[\tsc{insP}
    [\tsc{f4}]
    [\tsc{datP}
        [\tsc{datP}
            [\tsc{f3}]
            [\tsc{accP}
               [\tsc{f2}]
               [\tsc{nomP}
                   [\tsc{f1}]
                   [ΣP
                       [Σ]
  										 [IP
  										 		[I]
                           [\tsc{thingP}
                               [\tsc{thing}, roof]
                           ]
  											]
                   ]
               ]
            ]
        ]
        {\draw (.east) node[right]{⇒ \tit{t}}; }
        [NP
            [\tsc{pencil}, roof]
        ]
        {\draw (.east) node[right]{⇒ \tit{potlood}}; }
    ]
]
    \end{forest}

\end{frame}

\begin{frame}{Summary}

    \begin{enumerate}
      \item Q: why *\tit{met dat} (but \tit{daar mee})? \pause\\
      A: \tit{daar mee} is the default, \tit{met dat} only appears under special circumstances (morphological derivation, cf. \tit{mice} and not \tit{mouse-s})\pause
      \item Q: why does the adposition change form? \pause\\
      A: \tit{mee} and \tit{met} are stored as different forms in the lexicon \pause
      \item Q: why does the adposition change position? \pause\\
      A: \tit{mee} and \tit{met} are stored as preposition or postposition in the lexicon \pause
      \item Q: why is there an \tit{r}, even though there is no location involved? \pause\\
      A: \tit{r} is ambiguous between \tsc{thing} and \tsc{location}
    \end{enumerate}

\end{frame}


\end{document}
